\documentclass[11pt]{article}
\usepackage{icmlko}
\begin{document}

\title{안 본 곳에서도 위 절반은 가려낸다 --- 그러나 약한 것들끼리는 못 세운다}
\author{월드모델 실험실}
\date{노트 575 · 2026년 8월 5일}
\maketitle

\begin{abstract}
노트 $571$ 이 ``비게임 앱 전이가 팔 $19$ 개 중 $16$ 개에서 음수인데 KR 만화는
$6$ 개뿐''을 남겼다. 왜인지 캐는 데 노트가 넷 걸렸고, 답은 \textbf{두 짝을
라벨 크기로 가르는} 순간 나왔다.

\begin{center}
\begin{tabular}{lrrr}
\hline
 & 전체 & \textbf{아래 절반} & 위 절반 \\
\hline
모바일(학습) & $+0.6262$ & $+0.4300$ & $+0.3032$ \\
KR 만화(집 밖) & $+0.6374$ & $\mathbf{+0.0688}$ & $+0.3400$ \\
비게임 앱(집 밖) & $+0.2897$ & $\mathbf{+0.0759}$ & $+0.2137$ \\
\hline
\end{tabular}
\end{center}

\textbf{전체 $\rho$ 는 두 배 넘게 다른데 아래 절반은 거의 같다}
($+0.069$ 대 $+0.076$). 앱이 낮은 것은 전이가 안 되어서가 아니라
\textbf{표본이 아래에 치우쳐서}다(라벨 중앙이 모바일의 $1/4.4$).
\end{abstract}

\begin{figure}[t]
\centering
\includegraphics[width=\textwidth]{figs/f1.pdf}
\caption{\textbf{가} --- 아래 절반은 같다. \textbf{나} --- 손실의 자리.
\textbf{다} --- 네 번 고쳤다.}
\end{figure}

\section{손실이 어디 있나}

\begin{center}
\begin{tabular}{lr}
\hline
앱 $-$ 모바일 & \\
\hline
\textbf{위 절반} & $\mathbf{-0.0895}$ \\
\textbf{아래 절반} & $\mathbf{-0.3541}$ \\
\hline
\end{tabular}
\end{center}

\textbf{위 절반에서는 학습 도메인과 $0.09$ 차뿐이다.} 손실이 통째로 아래에
있다.

\begin{quote}
\textbf{KR 도 무너진다} --- 전체 $+0.6374$ 인데 아래 절반 $+0.0688$ 이다.
KR 의 높은 전체 $\rho$ 는 \emph{위·아래를 가르는 힘}에서 오고, 각 절반
안에서는 앱과 비슷하다. \textbf{거친 신호이지 고운 신호가 아니다.}
\end{quote}

\section{네 번 고쳐 답에 닿았다}

\begin{center}
\begin{tabular}{lll}
\hline
노트 & 읽기 & 판정 \\
\hline
$571$ & 팔 목록 정정 & $-0.34 \to -0.22\!\sim\!+0.10$ \\
$572$ & 앱은 다른 세계(축 뜻 $0.60$ 배) & \textbf{과장} \\
$573$ & 같은 세계의 아래쪽(모바일 꼬리) & \textbf{과장} \\
$574$ & $58\%$ 꼬리 $\cdot$ $42\%$ 전이 & 둘 다 절반씩 \\
\textbf{$575$} & \textbf{둘 다 아래 절반에서 무너진다} & \textbf{여기} \\
\hline
\end{tabular}
\end{center}

$572$ 는 ``축의 뜻이 학습과 $0.60$ 배''를 보고 다른 세계라 했고, $573$ 은
모바일 꼬리가 앱에 가까운 것을 보고 뒤집었다. \textbf{둘 다 짝을 통째로 놓고
잰 것이 문제였다} --- 반으로 가르니 ``세계가 다른가''가 아니라 ``어느
구간에서 되는가''가 물음이었다.

\section{제품 층에서 할 말}

\begin{quote}
\textbf{안 본 도메인에서 위 절반은 가려낸다}($+0.21 \sim +0.34$, 학습
도메인 위 절반 $+0.30$ 과 비슷하다). \textbf{약한 것들끼리는 못 세운다}
($+0.07$).
\end{quote}

노트 $536$(``쉬운 판만 샀다'')과 덱의 $1$ 대 $2$ 등 정확도 $57.4\%$(노트
$556$ · 격차가 $1.96$ 배뿐인 두 후보)가 \emph{집 안}에서 말한 것을, 이것이
\emph{집 밖}에서 말한다. \textbf{같은 이야기의 다른 판본이다} --- 가까운
것들끼리는 못 세운다.

\section*{거울이 왜 통과하기 어려운가}

노트 $426$ 거울(집 밖 짝 둘의 부호가 같아야 한다)은 팔 $19$ 개 중
\textbf{$3$ 개}만 통과했고 그 셋 중 둘이 채택으로 이어졌다(노트 $498 \cdot
553$). 이제 이유가 보인다 --- \textbf{앱 짝은 절반 이상이 ``못 세우는''
구간에 있으므로}, 그 짝을 올리려면 아래 절반을 건드려야 하는데 그것이 이
판에서 제일 어렵다.

\section*{조항}

\begin{quote}
\textbf{짝을 통째로 놓고 재면 ``세계가 다른가''로 읽게 된다.} 반으로 갈라
보면 ``어느 구간에서 되는가''가 나온다. \textbf{집 밖 성적을 한 숫자로
적지 않는다} --- 최소한 위/아래를 나눠 적는다.
\end{quote}

\section*{남은 것}

\begin{itemize}
\item \textbf{아래 절반이 왜 무너지나} --- 작은 수의 잡음인가 다른 기제인가.
라벨에 잡음을 모의로 넣어 위 절반 $\rho$ 가 얼마나 떨어지는지 재면 갈린다.
\item \textbf{노트 $443$ 의 \texttt{calibrate\_ooc}} --- 집 밖 축의 부호를
새 수집 라벨 몇 개로 맞추는 손잡이가 있는데 지금 안 쓴다. 아래 절반에
도움이 되나.
\item \textbf{셋째 집 밖 짝} --- 지금 둘뿐이고 둘 다 아래에서 무너진다.
셋째가 있으면 그것이 일반인지 이 둘의 성질인지 갈린다.
\end{itemize}

\end{document}
